% =====================================================================
% Khi nào phối hợp đa tác tử LLM tạo ra giá trị?
% Bản v1.1 — biên dịch: tectonic (XeTeX) hoặc pdflatex/Overleaf.
% v1.1: rà giọng văn trung tính theo phản biện độc lập.
% Ghi chú rà soát: \todoD{...} = Đức, \todoTD{...} = Tùng Dương
% =====================================================================
\documentclass[11pt,a4paper]{article}
% Tu nhan dien engine: XeTeX/LuaTeX (tectonic, xelatex) dung fontspec;
% pdfLaTeX (Overleaf mac dinh) dung T5 + inputenc. Ca hai deu ra dau tieng Viet dung.
\usepackage{iftex}
\ifPDFTeX
  \usepackage[T5]{fontenc}
  \usepackage[utf8]{inputenc}
  \usepackage[vietnamese]{babel}
\else
  \usepackage{fontspec}
  % Nap theo ten file de tuong thich bundle cua tectonic
  \setmainfont{texgyretermes}[Extension=.otf, UprightFont=*-regular,
    BoldFont=*-bold, ItalicFont=*-italic, BoldItalicFont=*-bolditalic]
  \setsansfont{texgyreheros}[Extension=.otf, UprightFont=*-regular,
    BoldFont=*-bold, ItalicFont=*-italic, BoldItalicFont=*-bolditalic]
  \setmonofont{texgyrecursor}[Extension=.otf, UprightFont=*-regular,
    BoldFont=*-bold, ItalicFont=*-italic, BoldItalicFont=*-bolditalic]
  \usepackage[vietnamese]{babel}
\fi
\usepackage{amsmath,amssymb}
\usepackage{booktabs}
\usepackage{multirow}
\usepackage[margin=2.4cm]{geometry}
\usepackage{microtype}
\usepackage{xcolor}
\usepackage[colorlinks=true,linkcolor=blue!60!black,citecolor=blue!60!black,urlcolor=blue!60!black]{hyperref}

\newcommand{\todoD}[1]{\textcolor{red!70!black}{[\textbf{Đ}: #1]}}
\newcommand{\todoTD}[1]{\textcolor{red!70!black}{[\textbf{TD}: #1]}}
\newcommand{\mucB}{\textsuperscript{\,(mức B)}}
\newcommand{\dceil}{\Delta_{\mathrm{ceil}}}
\newcommand{\dhonest}{\Delta_{\mathrm{honest}}}
\newcommand{\CEIL}{\mathrm{CEIL}}
\newcommand{\acc}{\mathrm{acc}}

\title{\bfseries Khi nào phối hợp đa tác tử LLM tạo ra giá trị?\\
\large Một khảo sát thực nghiệm dưới ba phép kiểm soát}
\author{Nguyên \and Đức \and Tùng Dương\thanks{Thứ tự tên và việc ghi nhận Quân: quyết ở Bước 0.
Bản v1.0: \S2 do Nguyên chấp bút tạm, Đức rà và xác nhận trích dẫn; \S4 Tùng Dương rà.}}
\date{Tháng 8, 2026}

\begin{document}
\maketitle

\begin{abstract}
Các hệ đa tác tử LLM --- planner--solver--verifier--aggregator, tranh biện, sinh-rồi-sửa ---
thường được báo cáo là cải thiện đáng kể so với một model đơn lẻ. Khảo sát này đo lại nhận định
đó trên GSM8K, MATH-500, MBPP và HumanEval với Qwen2.5 (1.5B--32B), Llama-3.1-8B và
DeepSeek-Coder-6.7B, dưới \emph{ba phép kiểm soát} mà phần lớn báo cáo trong lĩnh vực bỏ qua:
(1) \textbf{cùng ngân sách suy luận} --- pipeline bốn vai tốn 2,9--6,6 lần chi phí một lượt
giải, và token của các model khác cỡ không cùng đơn giá; (2) \textbf{đúng mốc so sánh} --- người
triển khai đã có model mạnh trong tay, nên đại lượng đúng là chênh lệch với \emph{model mạnh
chạy một mình} ($V-I$), không phải với model yếu ($V-S$); (3) \textbf{đúng mẫu số} --- 57\% số
câu không thể chịu tác động của bất kỳ cơ chế phối hợp nào, khiến hiệu ứng thật bị pha loãng
3,3 lần.

Sau khi áp dụng ba phép kiểm soát, dấu của nhiều hiệu ứng thay đổi một cách nhất quán. Hiệu ứng dương lớn nhất ($+14{,}0$ điểm, 5/5
fold) chỉ tồn tại khi so với model yếu; so cùng chi phí FLOP với model mạnh đơn lẻ thì hoà hoặc
kém. Nguyên nhân được định vị bằng một đẳng thức đại số $\dceil = A - B + C$ (kiểm tra nhất quán số liệu: khớp 19/19 cặp đo) và một thí nghiệm tiền đăng ký tái lập trên hai miền: thiệt hại đến từ việc model
mạnh \textbf{tiếp xúc với nội dung sai} --- trên MATH, độ chính xác của nó rơi từ 46,4\% xuống
19,2\% trên tầng artifact sai ($-27{,}2$ điểm; $p \approx 0$), trong khi artifact đúng cải thiện độ chính
xác của nó ($+3{,}8$ điểm; $p = 0{,}012$). Cùng một biến chênh lệch năng lực nuôi cả cơ hội
($A$) lẫn thiệt hại ($B$): giao thức \textbf{tuyển chọn} có $B \approx 0$ theo cấu trúc nên
chênh lệch càng lớn càng thắng; giao thức \textbf{sửa chữa} phải trả $B$ nên chênh lệch càng lớn
càng thua; điều này giải thích sự khác dấu biểu kiến giữa hai giao thức.

Ở miền có bộ kiểm \textbf{đúng đắn} (code chạy test được), phối hợp đạt đúng trần oracle và
không phá một bài nào trong 20/20 fold; ở miền không có, một bộ phân loại đạt AUC 0,893 vẫn chỉ
đổi được $+2{,}4$ điểm --- vì giá trị của mọi bộ kiểm bị chặn bởi khoảng cách
$\texttt{oracle@}k - \texttt{maj@}k$, không bởi chất lượng tín hiệu. Trong phạm vi khảo sát
(model 1,5B--32B, bốn benchmark suy luận), cấu hình sinh độc lập rồi tuyển chọn cho kết quả cao
hơn các cấu hình cho model mạnh sửa bài làm của model yếu.
\end{abstract}

% =====================================================================
\section{Mở đầu}\label{sec:intro}

Khảo sát này xét câu hỏi: ghép nhiều vai LLM có tạo thêm giá trị so với một model đơn lẻ
không, và nếu có, giá trị đến từ \emph{cấu trúc phối hợp} hay từ các yếu tố đi kèm cấu trúc đó.
Câu hỏi này khó trả lời hơn vẻ ngoài vì ba vấn đề đo lường.

\paragraph{Thứ nhất, ngân sách.} Pipeline bốn vai $P{\to}S{\to}V{\to}A$ tốn 2,9 lần ký tự sinh
trên GSM8K và 6,63 lần trên MATH so với một lượt giải (\S\ref{sec:start}). So một hệ chi phí
$3\times$ với một hệ $1\times$ rồi kết luận ``phối hợp có giá trị'' là so
sánh hai hệ ở hai mức ngân sách khác nhau, không tách được hiệu ứng phối hợp khỏi hiệu ứng chi
phí.
Tinh vi hơn: khi các vai dùng model \emph{khác cỡ}, ngay cả số token cũng không cùng đơn vị ---
một token do model 7B sinh tốn $\sim$4,7--4,9 lần FLOP so với token của model 1.5B, vì chi phí
suy luận tỷ lệ với số tham số \cite{kaplan2020} (\S\ref{sec:baseline}).

\paragraph{Thứ hai, mốc so sánh.} Gần như toàn bộ dòng sinh-rồi-sửa đo cải thiện so với model
\emph{bị sửa} --- tức model yếu ($V-S$). Nhưng lựa chọn thực của người triển khai, vốn đã có
model mạnh trong pipeline, là: chạy pipeline, hay gọi thẳng model mạnh ($V-I$)? Khảo
sát này cho thấy hai mốc đó cho hai dấu ngược nhau trên cùng một hệ
(\S\ref{sec:baseline}, \S\ref{sec:exposure}).

\paragraph{Thứ ba, mẫu số.} Phân tầng theo độ khó cho thấy 57\% số câu nằm ngoài tầm với của mọi
cơ chế phối hợp: 32\% quá sức (không mẫu nào đúng --- không có gì để chọn), 25\% quá dễ (mọi mẫu
đều đúng --- không có gì để cải thiện). Hiệu ứng $+31{,}1$ điểm trên tầng khả tác động bị pha
loãng thành $+9{,}3$ khi trung bình toàn tập (\S\ref{sec:denominator}). Thiếu kiểm soát này,
hiệu ứng thật bị đánh giá \emph{thấp} --- ngược chiều với hai kiểm soát trên.

Khảo sát hợp nhất ba nhánh nghiên cứu độc lập của nhóm --- phân rã đóng góp theo vai bằng giá
trị Shapley \cite{shapley1953}, phân tích chi phí--lợi ích của cơ chế định tuyến, và phân tích
luồng thông tin của giao thức sửa chữa --- dưới một khung duy nhất, và trình bày kết quả theo
mạch dẫn dắt: mỗi thí nghiệm sinh ra từ câu hỏi mà thí nghiệm trước để lại.

\paragraph{Đóng góp.}
\begin{enumerate}\itemsep2pt
  \item Khung phân rã giá trị $\text{value} = H(\text{pool}) \times \kappa(z) -
        D(\text{protocol})$ thống nhất ba nhánh (\S\ref{sec:framework}).
  \item Bằng chứng trên hai thiết kế độc lập rằng \textbf{chênh lệch năng lực} --- không phải
        số lượng vai, danh tính vai, hay họ model --- là biến giải thích chính trong các
        thí nghiệm của chúng tôi (\S\ref{sec:roles},
        \S\ref{sec:termA}).
  \item Một \textbf{đẳng thức đại số} $\dceil = A - B + C$ tách cơ hội khỏi thiệt hại; kiểm tra
        nhất quán số liệu khớp 19/19 cặp đo (\S\ref{sec:framework}, \S\ref{sec:termA}).
  \item Bằng chứng tiền đăng ký, tái lập hai miền: thiệt hại của giao thức sửa đến từ
        \textbf{tiếp xúc với nội dung sai}, không phải từ tiếp xúc; và giao thức sửa hoạt động như
        \textbf{ống dẫn} hơn là bộ sửa lỗi --- truyền đáp án đúng của model yếu (90,9\%, trên
        tầng $n=11$)
        nhưng gần như không sáng tạo lời giải mới (4,3\%) và phá ở quy mô lớn (63,6\% trên tầng
        rủi ro) (\S\ref{sec:exposure}).
  \item Giải nghịch lý ``cùng biến, hai dấu'' của chênh lệch năng lực (\S\ref{sec:synthesis}).
  \item Một phát biểu định lượng về $\kappa$: giá trị của bộ kiểm bằng khoảng cách
        $\texttt{oracle@}k - \texttt{maj@}k$, không bằng chất lượng bộ kiểm --- kèm hai trần
        định tuyến tính được (\S\ref{sec:kappa}).
  \item Một bộ quy trình chống tự đánh lừa (tiền đăng ký, điều kiện hợp lệ, niêm phong hash,
        sàn nhiễu, t-test theo fold); 50\% số lần chạy không vượt qua bộ điều kiện này ---
        chúng tôi báo cáo tỷ lệ đó như một quan sát về quy trình (\S\ref{sec:method}).
\end{enumerate}

% =====================================================================
\section{Công trình liên quan}\label{sec:related}

\paragraph{Sinh-rồi-chọn.} Chuỗi suy luận \cite{wei2022} mở đường cho self-consistency
\cite{wang2023}: sinh $k$ lời giải độc lập rồi bỏ phiếu đa số --- chính là \texttt{maj@}$k$ của
khảo sát này. Dòng chọn-bằng-bộ-kiểm gồm chấm theo quá trình \cite{lightman2024} và mở rộng
compute lúc suy luận \cite{snell2024}. Đặc điểm chung: giao thức \emph{không cho phép tạo lời
giải mới ở bước tổng hợp}, tức $B \approx 0$ theo thuật ngữ \S\ref{sec:identity} của chúng tôi.

\paragraph{Sinh-rồi-sửa.} Self-Refine \cite{madaan2023}, Reflexion \cite{shinn2023} và CRITIC
\cite{gou2024} cho model đọc sản phẩm (của chính nó hoặc model khác) rồi sửa. Huang và cộng sự
\cite{huang2024} chỉ ra LLM chưa tự sửa được suy luận khi thiếu tín hiệu ngoài --- nhất quán với
kết quả \texttt{llm3} và phân tầng tiếp xúc của chúng tôi. Theo rà soát sơ bộ của chúng tôi, phần lớn
công trình dòng này báo cáo cải thiện so với model bị sửa hoặc lượt nháp đầu ($V-S$), không
so với việc dùng thẳng model sửa ($V-I$) \todoD{rà từng bài, xác nhận mốc so sánh của
từng bài --- đây là chỗ luận điểm \S\ref{sec:intro} đứng hoặc đổ}.

\paragraph{Đa tác tử và tranh biện.} Tranh biện đa tác tử \cite{du2024} và LLM-làm-giám-khảo
\cite{zheng2023} là các cơ chế phối hợp phổ biến. Tổng hợp số liệu công bố gần đây của nhóm
\todoD{nguồn cụ thể} cho thấy tranh biện thua self-consistency ở đa số ô so sánh và giảm
16 điểm với Llama-3.1-8B trên GSM8K ($80{,}13 \to 63{,}87$) --- nhất quán với mẫu hình suy giảm ở
model nhỏ mà chúng tôi quan sát độc lập (\S\ref{sec:roles}). Hệ gần nhất về kiến trúc là
MAS\_RPSV (bốn vai nối tiếp, cùng cỡ model, cùng benchmark) và SHARP --- cùng dùng giá trị
Shapley cho credit assignment, và báo cáo chỉ 12,96\% lượt gọi subagent là hữu ích \todoD{xác
nhận cả hai nguồn} --- song song với kết quả ``2/4 vai thực sự tính toán'' của chúng tôi.

\paragraph{Vị trí của khảo sát.} Chúng tôi không đề xuất kiến trúc mới; chúng tôi đề xuất
\textbf{cách đo}: ba phép kiểm soát, một đẳng thức phân rã, và một khung ba thừa số --- rồi dùng
chúng để trả lời \emph{khi nào} các kiến trúc trên tạo giá trị thật.

% =====================================================================
\section{Khung lý thuyết}\label{sec:framework}

\subsection{Phân rã giá trị}
\begin{equation}
\text{value} \;=\; H(\text{pool}) \times \kappa(z) \;-\; D(\text{protocol})
\end{equation}
\begin{itemize}\itemsep1pt
  \item $H$ --- \textbf{dư địa}: pool ứng viên có chứa lời giải mà model mạnh đơn lẻ không tạo
        ra được không? Đo bằng \texttt{oracle@}$k$.
  \item $\kappa$ --- \textbf{chất lượng bộ chọn}: một tín hiệu \emph{khả thi} (không phải
        oracle) có lấy được lời giải đó ra khỏi pool không?
  \item $D$ --- \textbf{thiệt hại}: bản thân giao thức phá hỏng bao nhiêu?
\end{itemize}
Ba nhánh của nhóm đo ba thừa số: Shapley theo vai đo $H$; router đo $\kappa$; luồng thông tin
đo $D$.

\subsection{Đẳng thức phân rã}\label{sec:identity}

Ký hiệu (chú ý: \emph{không} trùng với tên vai ở \S\ref{sec:setup}): $S$ = model yếu giải bài;
$I$ = model mạnh giải bài, \emph{không thấy gì của} $S$; $V$ = \textbf{cùng model mạnh đó}
nhưng ngữ cảnh có kèm bài làm của $S$. $I$ và $V$ là cùng một model --- khác nhau ở
\textbf{input}, nên $V-I$ cô lập đúng một biến: việc nhìn thấy bài làm của model yếu. Trần lý
tưởng $\CEIL$ = đúng nếu $S$ đúng \textbf{hoặc} $V$ đúng --- chính là \texttt{oracle@2} trên cặp
$\{S, V\}$, tức chiến lược ``$S$ đúng thì giữ $S$'' ở dạng hoàn hảo.
\begin{equation}
\dceil \;=\; \acc(\CEIL) - \acc(I) \;=\; A - B + C
\end{equation}
\begin{align*}
A &= P(S \text{ đúng} \wedge I \text{ sai})   &&\text{cơ hội có sẵn --- tính chất của \emph{cặp model}}\\
B &= P(S \text{ sai} \wedge I \text{ đúng} \wedge V \text{ sai}) &&\text{giao thức phá bài $I$ vốn đúng --- tính chất của \emph{giao thức}}\\
C &= P(S \text{ sai} \wedge I \text{ sai} \wedge V \text{ đúng}) &&\text{giao thức cứu bài cả hai đều sai}
\end{align*}
Đây là đẳng thức đại số, không phải mô hình: khớp tuyệt đối trên 4/4 cặp có lưu vết và 15/15 cặp
của thí nghiệm quy mô lớn. Giao thức \textbf{chỉ chọn} có $B \approx 0$ vì không thể tạo lời
giải sai mới; giao thức \textbf{sửa} luôn phải trả $B$. $\dceil$ là \emph{chặn trên sàng lọc}
(cần đáp án chuẩn, không triển khai được): $\dceil < 0 \Rightarrow$ đóng hướng ngay --- không
cổng nào cứu nổi; $\dceil > 0 \Rightarrow$ mới được hỏi tiếp về $\kappa$. Con số triển khai là
$\dhonest = \acc(V) - \acc(I)$.

\subsection{Ba mệnh đề}
\begin{itemize}\itemsep1pt
  \item[\textbf{M1.}] $D$ là hàm của tiếp xúc với \textbf{nội dung sai}, không phải của tiếp
        xúc (\S\ref{sec:exposure}).
  \item[\textbf{M2.}] $\kappa$ phụ thuộc tín hiệu là \textbf{đúng đắn} hay \textbf{học được}
        --- không phụ thuộc mạnh hay độc lập. (Đã qua hai lần thu hẹp từ bản gốc ``độc lập
        thắng tương quan''; \S\ref{sec:kappa}.)
  \item[\textbf{M3.}] Định tuyến tiết kiệm ở đúng nơi ít cần tiết kiệm nhất (\S\ref{sec:router}).
\end{itemize}

% =====================================================================
\section{Thiết lập thí nghiệm}\label{sec:setup}

\todoTD{rà toàn bộ mục này}

\paragraph{Model.} Qwen2.5-Instruct 1.5B/7B/14B/32B \cite{yang2024}; Llama-3.1-8B-Instruct
\cite{grattafiori2024}; DeepSeek-Coder-6.7B-Instruct \cite{guo2024}. Lượng tử hoá \textbf{đo
từng model, không suy theo họ}: DeepSeek-6.7B nạp nf4 hết 3,61 GB; Qwen-7B 5,21 GB; Llama-8B và
Qwen-14B \emph{không} lượng tử hoá được trên bản chạy (fp16 $\sim$16/28 GB). Phần cứng: Kaggle
2$\times$T4 (nf4/fp16) và RTX 6000 Pro 95 GB (bf16); so sánh chéo lần chạy chỉ hợp lệ khi trùng
cả (phần cứng + độ chính xác số) lẫn bộ bài --- đo được nf4 so với bf16 dịch kết quả tới 0,03.

\paragraph{Benchmark.} GSM8K \cite{cobbe2021}; MATH-500 \cite{hendrycks2021}; MBPP
\cite{austin2021} --- dải chính \texttt{task\_id} 11--510 ($\sim$499 bài sau lọc chuẩn), dải giữ
lại 511--974 (464 bài); HumanEval \cite{chen2021}. \textbf{MBPP/HumanEval có bộ kiểm đúng đắn}
(chạy test); MATH/GSM8K không --- khác biệt này quyết định kết luận \S\ref{sec:kappa}.

\paragraph{Vai và chi phí.} $P$ planner, $S$ solver, $V$ verifier, $A$ aggregator; tổ hợp
\texttt{PS}, \texttt{PSV}, \texttt{PSVA}, \texttt{SVV},\dots{} Chi phí báo theo \textbf{lượt
gọi} và theo \textbf{token/ký tự}; khi các vai khác cỡ model, quy thêm về \textbf{chi phí trọng
số FLOP} (tham số $\times$ token).

\paragraph{Giải mã và tính tất định.} Khối luồng thông tin dùng greedy
(\texttt{do\_sample=False}): cùng cấu hình cho kết quả \textbf{giống hệt từng bài} (kiểm 499/499
giữa hai tài khoản, hai ngày). Hệ quả: chạy lại cùng cấu hình không phải bằng chứng độc lập; dao
động giữa các fold đến từ \emph{khác bài}, không phải nhiễu lấy mẫu.

\paragraph{Thống kê.} Khối vai trò/chi phí: 5 fold rời nhau; sàn nhiễu hiệu chuẩn bằng cách chạy
cùng cấu hình trên 5 fold (\texttt{V\_gain} dao động $+1{,}0\ldots+8{,}0$; std giữa fold 2,65);
suy diễn bằng t-test trên fold ($\mathrm{SE} = s/\sqrt{k}$; $t_{0{,}975}(4) = 2{,}776$), khoảng tin cậy 95\% (KTC),
và phép thử dấu (5/5 cùng dấu $\Rightarrow p = 0{,}031$ một phía). Khối luồng thông tin:
McNemar chính xác ghép cặp trên cùng bộ bài; bootstrap ghép cặp theo chỉ số bài (10\,000 lần).
Ngưỡng diễn giải được \textbf{khoá trước khi chạy} trong bảng tiền đăng ký, gồm dòng bác bỏ giả
thuyết.

\paragraph{Nhãn tin cậy và đơn vị.} \emph{Mức A} = tiền đăng ký với điều kiện hợp lệ, hoặc 5
fold kèm t-test; \emph{mức B} = đo một lần hoặc phân rã hậu nghiệm --- chỉ dùng minh hoạ cơ
chế, không dùng làm kết luận chính. Accuracy viết ở thang 0--1; ``điểm'' nghĩa là điểm phần
trăm ($+14{,}0$ điểm $= +0{,}140$).

% =====================================================================
\section{Hành trình thực nghiệm}\label{sec:results}

\emph{Trình bày theo mạch: mỗi mục nối tiếp vấn đề mà mục trước để lại. Số liệu ở mức A trừ
khi ghi khác (định nghĩa nhãn: \S\ref{sec:setup}).}

% ---------------------------------------------------------------
\subsection{Pipeline có giúp không --- và sàn nhiễu định hình mọi thứ}\label{sec:start}

\begin{table}[h]\centering\small
\begin{tabular}{lcccc}\toprule
Task & Solver & Pipeline & Chênh & Chi phí ký tự\\\midrule
GSM8K & 0,632 & 0,744 & $\mathbf{+0{,}112}$ & $2{,}9\times$\\
MATH  & 0,405 & 0,345 & $\mathbf{-0{,}060}$ & $\mathbf{6{,}63\times}$\\\bottomrule
\end{tabular}
\caption{Pipeline đầy đủ so với solver đơn lẻ (Qwen-1.5B).}
\end{table}

Trên GSM8K có lợi (bảng trên so pipeline với solver một lượt trên toàn tập; kiểm định thống
kê dùng đại lượng 5 fold $\texttt{PSVA}-\texttt{PS} = +5{,}6$; KTC $[+3{,}3;\,+7{,}9]$;
$t=6{,}89$); trên MATH chi phí gấp 6,63 lần nhưng độ chính xác thấp hơn 6 điểm. Phép hiệu chuẩn sàn nhiễu ---
cùng cấu hình, 5 fold --- cho \texttt{V\_gain} từ $+1{,}0$ đến $+8{,}0$: một phép đo đơn lẻ 100
bài gần như không mang thông tin, và mọi hiệu ứng từ đây đi kèm KTC.

\emph{Mục tiếp theo kiểm tra phần lợi này đến từ phối hợp hay từ số lượt sinh.}

% ---------------------------------------------------------------
\subsection{Kiểm soát ngân sách: giá trị nằm ở lượt sinh, không ở vai tổng hợp}\label{sec:budget}

\begin{table}[h]\centering\small
\begin{tabular}{lcc}\toprule
Bộ tổng hợp (8 lượt sinh) & MATH 1.5B & MATH 7B\\\midrule
greedy (1 lượt)                 & 0,50 & 0,72\\
\texttt{maj@8} --- bỏ phiếu     & \textbf{0,60} & \textbf{0,73}\\
\texttt{llm\_agg@8} --- LLM tổng hợp & 0,41 & 0,47\\
\texttt{oracle@8} --- trần      & 0,73 & 0,85\\\bottomrule
\end{tabular}
\caption{Cùng ngân sách 8 lượt sinh, chỉ đổi bộ tổng hợp.}
\end{table}

Cùng chi phí, thành phần LLM tổng hợp kém bỏ phiếu 19--26 điểm; ở 7B kém cả greedy một lượt. Nó
phá đa số đúng 21/26 lần (1.5B/7B), sửa được đa số sai 2/0 lần. Kết quả phụ thuộc task
($\texttt{vote5}-\texttt{llm\_agg}$: $+0{,}075$ MATH, $-0{,}056$ GSM8K); phát biểu phù hợp với
dữ liệu: \emph{bộ tổng hợp LLM không mang giá trị ổn định giữa các task; trên MATH nó kém bỏ
phiếu 19--26 điểm.}

Hai đối chứng tách nốt:
\begin{itemize}\itemsep2pt
\item \textbf{\texttt{rerun} = \texttt{loop}} ($0{,}453 = 0{,}453$): giải lại \emph{không xem}
      phê bình của verifier cho kết quả bằng chính xác giải lại \emph{có} phê bình. Lợi ích đến
      từ \textbf{lượt sinh thêm}, không từ nội dung phản hồi. (Một phép đo đơn lẻ trước đó cho $+20$
      nhưng không tái lập: trên 5 fold còn $+4{,}0$.) Biến thể vòng lặp giải-rồi-chấm rẻ hơn pipeline đủ vai lại
      \textbf{thắng nó} trên MATH: $\texttt{loop}-\texttt{PSVA} = +4{,}0$; KTC $[+0{,}5;\,+7{,}5]$;
      $t = 3{,}21$ --- cùng chi phí (4,2 so với 4,0 lượt).
\item \textbf{Aggregator chủ yếu sao chép thay vì lựa chọn:} 65\% số câu nó chép nguyên ứng viên
      \emph{cuối cùng}; đọc 600 trace: 100\% lượt không sinh số mới. \textbf{Nhưng phần ``gây
      hại'' phải đính chính:} 85\% số ca phá trên MATH là \emph{không phát ra}
      \verb|\boxed{}| --- lỗi định dạng; chỉ 5\% chọn nhầm thật. Thêm fallback miễn phí:
      \texttt{A\_gain} từ $-6{,}4$ thành $\mathbf{+1{,}0}$ $[0;\,+2]$. Bộ tổng hợp LLM
      \emph{trung tính một khi xử lý xong định dạng} --- không cộng giá trị, nhưng ``phá hoại''
      phần lớn là hiện vật đo lường. (Một mâu thuẫn chưa giải giữa hai lần chạy 5 fold của chính
      đại lượng này --- $-6{,}4$ và $+8{,}0$, mỗi bên nhất quán nội bộ --- xem \S\ref{sec:limits}.)
\end{itemize}

\emph{Các can thiệp trên đều cho 0--5 điểm; mục tiếp theo kiểm tra mẫu số của phép đo.}

% ---------------------------------------------------------------
\subsection{Kiểm soát mẫu số: 57\% số câu bất động}\label{sec:denominator}

\begin{table}[h]\centering\small
\begin{tabular}{lcccc}\toprule
Mẫu đúng /5 & $n$ & Solver & \texttt{vote5} & $\Delta$\\\midrule
0/5 --- quá sức & 48 (32\%) & 0,000 & 0,000 & 0,000\\
1/5 & 20 & 0,000 & 0,000 & 0,000\\
2/5 & 15 & 0,333 & 0,600 & $+0{,}267$\\
3/5 & 12 & 0,583 & 1,000 & $+0{,}417$\\
4/5 & 18 & 0,722 & 1,000 & $+0{,}278$\\
5/5 --- quá dễ & 37 (25\%) & 1,000 & 1,000 & 0,000\\\bottomrule
\end{tabular}
\caption{Phân tầng theo số mẫu đúng trong 5 lượt (MATH 1.5B, $n=150$).}
\end{table}

Với 57\% số câu, mọi bộ tổng hợp trên cùng pool đều cho cùng kết quả --- hệ quả cấu trúc
của phép chọn. Trên tầng khả tác động: $+21{,}5$ điểm (nhóm 1--4/5) hay $+31{,}1$ (nhóm 2--4/5, loại tầng hiệu ứng
bằng 0); trung bình toàn tập pha loãng còn $+9{,}3$. Một can thiệp $+10$ điểm thật sẽ hiện ra
$+3$ --- dưới sàn nhiễu, và có nguy cơ bị đánh giá nhầm là không có hiệu ứng. Ba kiểm soát tác động \textbf{ngược
chiều}: thiếu mẫu số $\Rightarrow$ đánh giá thấp; thiếu ngân sách/mốc $\Rightarrow$ đánh giá cao.

\emph{Mục tiếp theo phân rã đóng góp theo vai trên phần khả tác động.}

% ---------------------------------------------------------------
\subsection{Vai trò: chuyên biệt hoá không tạo hiệu ứng đo được; biến giải thích là năng lực của lượt kiểm}\label{sec:roles}

Chỉ số \textbf{hành vi} từ trace: ở 1.5B, planner --- bị cấm tính đáp án --- chứa sẵn đáp án ở
34,7\% câu MATH (45,3\% có \verb|\boxed|); solver không sinh số mới ở 60,7--62\% lượt, lời giải
dài 19 ký tự khi có plan (664 khi không) --- solver chủ yếu sao chép đáp án có sẵn trong plan
thay vì tự giải. Ở 7B, hai
vai hồi phục (rò đáp án 1,0--4,0\%) --- nhưng ở mức đó pipeline lại thua solver đơn lẻ. Đối
chứng nhân quả: \textbf{hoán vị prompt giữa các vị trí} (\texttt{normal}/\texttt{swap}/%
\texttt{solo}) cho accuracy không phân biệt được --- GSM8K $0{,}660 = 0{,}660$, \texttt{solo}
$0{,}673$; MATH $\texttt{swap}-\texttt{normal} = +5{,}3$ nhưng $t = 1{,}84$, KTC
$[-2{,}7;\,+13{,}4]$ --- không đạt ý nghĩa thống kê. Chúng tôi không tìm thấy bằng chứng
danh tính vai là biến có tác động; riêng phép thử MATH chưa đủ lực để loại trừ hiệu ứng cỡ vừa.

Biến giải thích là \textbf{năng lực của model ở lượt thứ hai}. Cùng 300 bài MATH (5 fold $\times$ 60):

\begin{table}[h]\centering\small
\begin{tabular}{lccc}\toprule
Lượt kiểm & Hiệu ứng & KTC 95\% & Sửa : Phá\\\midrule
\textbf{Verifier 7B} (solver 1.5B) & $\mathbf{+14{,}0}$ & $[+7{,}4;\,+20{,}6]$ & $\mathbf{43:1}$\\
Verifier 1.5B (cùng cỡ) & $+3{,}0$ & chạm 0 & $15:6$\\
Phần riêng do cỡ model lớn hơn & $\mathbf{+11{,}0}$ & $[+3{,}9;\,+18{,}1]$ & ---\\\bottomrule
\end{tabular}
\caption{Bất đối xứng năng lực ở lượt kiểm, MATH.}
\end{table}

Cơ chế từ trace: khi can thiệp, verifier tái sử dụng 0\% số liệu của solver --- nó
không dò từng bước mà giải lại từ đầu. Về hành vi, verifier là một lượt giải thứ hai chứ không
phải một bộ kiểm từng bước; lợi ích của verifier 7B tương đương lợi ích của việc thêm một lượt
giải bằng model mạnh hơn. Phù hợp: \texttt{V\_gain} có
ý nghĩa ở cả MATH 7B ($+4{,}4$; KTC $[+1{,}5;\,+7{,}3]$) --- giá trị verifier xuất hiện
\emph{cùng} năng lực; ở MATH 1.5B thì không ($+1{,}4$; $p=0{,}235$).

\emph{Con số $+14{,}0$ so với model yếu; mục tiếp theo so lại theo mốc $V-I$ và đơn giá
token.}

% ---------------------------------------------------------------
\subsection{Kiểm soát mốc so sánh và đơn giá token: cấu hình bất đối xứng bị chi phối}\label{sec:baseline}

\begin{table}[h]\centering\small
\begin{tabular}{lcccc}\toprule
Cấu hình & GSM8K & MATH & Token GSM8K & Token MATH\\\midrule
S1.5B + V7B & 0,810 & 0,563 & 105k & 119k\\
\textbf{S7B một mình} & \textbf{0,910} & \textbf{0,593} & 120k & 152k\\
S7B + V7B & 0,900 & \textbf{0,670} & 205k & 261k\\\bottomrule
\end{tabular}
\caption{Token đo thật, 5 fold. \emph{Cột token không so được trực tiếp giữa hai cỡ model ---
xem hiệu chỉnh bên dưới.}}
\end{table}

So với S7B đơn lẻ: bất đối xứng kém 10 điểm trên GSM8K; hoà trên MATH (chênh theo fold chứa 0).
Cột token có \textbf{hai lỗi chồng nhau}: (a) trường đo gốc chỉ đếm token \emph{do 7B sinh}, bỏ
sót token của solver 1.5B; (b) token 7B đắt $\sim$4,7--4,9 lần token 1.5B. Quy về trọng số FLOP
(tham số $\times$ token; độ nhạy chars/token 3,0--4,0):

\begin{table}[h]\centering\small
\begin{tabular}{lcc}\toprule
 & Token thô 7B & \textbf{Trọng số FLOP}\\\midrule
GSM8K: asym / S7B & 0,88 --- ``rẻ hơn 12\%'' & $\mathbf{1{,}00\text{--}1{,}05}$ --- tiết kiệm biến mất\\
MATH: asym / S7B & 0,78 --- ``rẻ hơn 22\%'' & $\mathbf{0{,}92\text{--}0{,}96}$ --- còn 3--8\%\\\bottomrule
\end{tabular}
\caption{Hiệu chỉnh đơn giá token. Chưa tính prefill (7B còn phải \emph{đọc} bài của 1.5B) nên
hiệu chỉnh vẫn thiên vị cho cấu hình bất đối xứng.}
\end{table}

Kết luận: trên GSM8K cấu hình bất đối xứng bị chi phối (kém 10 điểm với chi phí FLOP tương
đương); trên MATH là tiết kiệm \emph{biên} với độ chính xác hoà. Solver đơn lẻ nằm trên đường
Pareto ở cả hai task; pipeline đầy đủ không nằm trên đường Pareto ở MATH. Cải thiện duy nhất có ý
nghĩa thống kê trong bảng: thêm V7B vào \emph{chính} S7B trên MATH: $+7{,}7$ ($p = 0{,}0075$) --- với giá
$1{,}7\times$ token.

Đây là dạng tổng quát của vấn đề mốc: $V-S$ dương, $V-I$ âm --- cùng hệ, hai mốc, hai
kết luận ngược nhau --- lặp lại ở các nhánh của khảo sát.

\emph{Các mục tiếp theo dùng đẳng thức \S\ref{sec:identity} để định vị phần giá trị mất
đi.}

% ---------------------------------------------------------------
\subsection{Định tuyến: tiết kiệm ở nơi ít cần nhất \normalfont(mức B --- đo một lần, $n=300$)}\label{sec:router}

Consensus router (chạy $S$, $V$; bất đồng mới gọi $A$): GSM8K đạt 0,7200 với chi phí 2,32 lượt
--- giữ 94\% phần lợi của pipeline (0,7233; 3 lượt) với 77\% chi phí; MATH đạt 0,4133 ---
bằng đúng solver đơn lẻ, tức không mang lại cải thiện. Cơ chế: khi $S$--$V$ bất đồng, $A$ cứu 45,4\% trường
hợp trên GSM8K nhưng 25,0\% trên MATH. Mệnh đề M3: định tuyến chỉ trả tiền ở nơi đã sẵn đồng
thuận. (Vì sao MATH tệ hơn --- trả lời ở \S\ref{sec:exposure}: model yếu sai nhiều hơn
$\Rightarrow$ tiếp xúc nội dung sai nhiều hơn.)

% ---------------------------------------------------------------
\subsection{Số hạng $A$: chênh lệch năng lực, không phải họ model --- và quy luật của $\dceil$}\label{sec:termA}

Quan sát thô trên 7 cặp gộp nhiều lần chạy: cặp \emph{khác họ} có dư địa $A$ cao gần gấp đôi
($0{,}0597$ so với $0{,}0481$) --- nhưng các cặp khác họ cũng có chênh lệch năng lực nhỏ hơn
($0{,}130$ so với $0{,}167$); hai biến trộn hoàn toàn. Thiết kế tách (tiền đăng ký):
\textbf{6 model, 15 cặp có hướng, cùng 499 bài MBPP, một lần chạy}, bổ sung đúng các ô còn thiếu:
\begin{align*}
A &= \beta_0 + \beta_1\cdot\text{chênh} + \beta_2\cdot\text{khác\_họ}\\
\beta_1 &= -0{,}1922\ (p \approx 0) \qquad
\beta_2 = +0{,}0045,\ \text{KTC 95\% } [-0{,}0051;\,+0{,}0140],\ p = 0{,}31
\end{align*}
$R^2$ chỉ với chênh lệch $= 0{,}824$. KTC của $\beta_2$ nằm trọn dưới ngưỡng $+0{,}02$
đã khoá trước --- một kết quả null có thông tin: ``khác họ'' là tương quan giả; biến giải thích
là chênh lệch năng lực. (Kết quả đa dạng pool --- 3 model cho 2,70/3 ứng viên phân biệt so với 1,91/3 khi
lấy mẫu cùng model; 36,2\% $\to$ 6,5\% số bài chỉ một ứng viên --- được ghi là ``khác
\textbf{model}'': đối chứng khác-model-cùng-họ chưa chạy.)

Thêm nhánh sửa cho cả 15 cặp (tiền đăng ký) cho quy luật:
\begin{equation}
\dceil = +0{,}0218 - 0{,}2392 \times \text{chênh};\quad R^2 = 0{,}60;\ p = 10^{-5};\
g^\ast \approx 0{,}09
\end{equation}
Đẳng thức $A-B+C$ khớp 15/15. Đọc kèm bắt buộc: 0/15 cặp dương có ý nghĩa; 3/15 âm có ý nghĩa
$\Rightarrow$ chỉ phát biểu được chiều phủ định: khi chênh lệch vượt $\sim$0,09, giao thức sửa
cho hiệu ứng âm trên MBPP. Xác lập vùng dương cần $\sim$8$\times$ toàn bộ MBPP --- bất khả thi
trên benchmark này. Chênh giải thích 82\% phương sai của $A$ nhưng chỉ 35\% của $B$. \emph{(Ghi
chú cơ chế, mức B: ngưỡng bảo toàn mà $V$ phải vượt tăng theo chênh với hệ số 2,18 --- nhanh gấp
đôi khả năng bảo toàn thực của nó, hệ số 1,10; chênh lệch hệ số $p = 0{,}0066$. Hiệu ứng âm của giao
thức sửa phát sinh từ cấu trúc ngân sách này, không phải từ suy giảm năng lực nội tại của model
mạnh.)}

\emph{Mục tiếp theo xác định nguồn gốc của $B$.}

% ---------------------------------------------------------------
\subsection{Số hạng $B$: hình phạt của nội dung sai --- và giao thức sửa là ống dẫn}\label{sec:exposure}

Thiết kế sạch nhất (tiền đăng ký; MATH-500 để \textbf{đổi miền} so với quan sát thăm dò gốc trên
MBPP): hai nhánh dùng \textbf{cùng một lệnh giải}, khác duy nhất việc ngữ cảnh có kèm bài làm
của model yếu; phân tầng theo \textbf{nội dung} bài làm đó:

\begin{table}[h]\centering\small
\begin{tabular}{lccc}\toprule
 & MBPP 11--510\mucB & MBPP 511--974\mucB & \textbf{MATH-500 (mức A)}\\\midrule
Artifact \textbf{sai}  & $-0{,}1900$ & $-0{,}1927$ & $\mathbf{-0{,}2720}$ ($p\approx 0$; $n=261$)\\
Artifact \textbf{đúng} & $+0{,}0636$ & $+0{,}0245$ & $\mathbf{+0{,}0377}$ ($p=0{,}012$; $n=239$)\\\bottomrule
\end{tabular}
\caption{Hiệu ứng của việc nhìn thấy artifact, phân tầng theo nội dung (thí nghiệm chính của
khảo sát; hai cột MBPP mức B, cột MATH mức A).}
\end{table}

Trên tầng artifact sai, model mạnh rơi từ 46,4\% xuống 19,2\% --- trên chính những bài nó
vốn giải được gần một nửa; artifact đúng cải thiện độ chính xác của nó. Tổng hợp trọng số tái tạo \emph{chính
xác} $V-I = -0{,}1240$. $D$ không phải hình phạt của việc nhìn thấy artifact nói chung, mà của việc nhìn thấy nội
dung sai; và vì thiết kế không có lệnh sửa, chỉ riêng sự hiện diện của nội dung sai trong ngữ
cảnh đã đủ tạo hiệu ứng âm.

Bảng $2\times2\times2$ đầy đủ (chấm lại từ trace) trả lời câu hỏi tự nhiên \emph{``khi $I$ sai
thì sửa có giúp không?''}:

\begin{table}[h]\centering\small
\begin{tabular}{lcc}\toprule
Tầng & $n$ & $V$ đúng\\\midrule
$S$ đúng, $I$ đúng & 228 & 99,6\%\\
\textbf{$S$ đúng, $I$ sai} & \textbf{11} & \textbf{90,9\%}\\
$S$ sai, $I$ đúng \emph{(vùng rủi ro)} & \textbf{121} & 36,4\% --- tức phá 63,6\%\\
Cả hai sai & 140 & \textbf{4,3\%}\\\bottomrule
\end{tabular}
\caption{Bảng $2\times2\times2$ trên MATH ($n=500$).}
\end{table}

Ba đọc số: (1) \textbf{giao thức sửa là ống dẫn, không phải bộ sửa lỗi} --- truyền đáp án đúng
của $S$ ở 10/11 bài tầng tương ứng (90,9\%; tầng nhỏ, độ bất định lớn), gần như không sáng tạo
lời giải mới (4,3\%), phá ở quy mô lớn (63,6\%); (2) tầng khai thác được nhỏ hơn tầng rủi ro \textbf{11 lần} (11 so với 121 bài); (3)
``giữ nguyên $S$ khi $S$ đúng'' chính là oracle $\CEIL$ --- và ngay cả khi cung cấp oracle đó với chi
phí bằng không, mức tăng cũng chỉ là $+0{,}004$ trên MATH (đúng 2 bài/500), vì $V$ vốn đã giữ được 99,6\%/90,9\%; kể cả vậy
$\CEIL = 0{,}578$ vẫn thua $I = 0{,}698$. Không cấu hình sửa nào đã đo có $V > I$ (0/15 cặp
MBPP, tốt nhất $-0{,}030$; các cặp MATH ở \S\ref{sec:transfer}). Vấn đề không phải ``quên giữ
$S$'' --- mà là $V$ phá tầng ($S$ sai, $I$ đúng); trong các phương án đã thử, chỉ việc không
đưa artifact vào ngữ cảnh bảo vệ được tầng này.

Đây cũng là lời giải cho \S\ref{sec:router}: aggregator vô dụng trên MATH vì ở đó model yếu sai
nhiều hơn --- tỷ lệ tiếp xúc nội dung sai cao hơn (45,4\% so với 25,0\% khớp hướng; mức B).

% ---------------------------------------------------------------
\subsection{Tính chuyển miền của quy luật chênh lệch}\label{sec:transfer}

Dùng đường khớp \textbf{trên MBPP} dự báo $\dceil$ \textbf{trên MATH} (tiền đăng ký; 3 cặp Qwen;
mọi điều kiện hợp lệ đạt; KTC bootstrap ghép cặp 10\,000 lần):

\begin{table}[h]\centering\small
\begin{tabular}{lccccl}\toprule
Cặp & Chênh & Đo được & KTC 95\% & MBPP dự báo & \\\midrule
7B$\to$14B & 0,044 & $-0{,}0140$ & $[-0{,}046;\,+0{,}018]$ & $+0{,}0108$ & trong khoảng\\
1.5B$\to$7B & 0,244 & $\mathbf{-0{,}1660}$ & $[-0{,}208;\,-0{,}124]$ & $-0{,}0361$ & \textbf{ngoài}\\
1.5B$\to$14B & 0,288 & $-0{,}0680$ & $[-0{,}102;\,-0{,}034]$ & $-0{,}0471$ & trong khoảng\\\bottomrule
\end{tabular}
\caption{Thử điểm tính chuyển miền của quy luật chênh lệch.}
\end{table}

2/3 dự báo trong khoảng $\Rightarrow$ quy luật \textbf{không bị bác bỏ} trên miền toán (không
viết ``đã xác nhận'': KTC rộng 0,064--0,084 nên phép kiểm có độ phân giải thấp). Cặp nằm ngoài khoảng lệch một
cách hệ thống --- $B = 0{,}208$, gấp $\sim$11 lần $A$ --- và tái lập một phép đo độc lập trên
cùng cặp ($-0{,}1380 \to -0{,}1660$): quy luật mô tả xu hướng; ngoại lệ ở nơi model yếu quá yếu
so với bài.

% ---------------------------------------------------------------
\subsection{$\kappa$: tín hiệu đúng đắn đạt trần oracle; tín hiệu học được không chuyển thành độ chính xác}\label{sec:kappa}

\textbf{Tín hiệu đúng đắn} --- HumanEval, cùng model, cùng 4 lượt sinh, khác duy nhất nguồn
kiểm; tái lập trên 4 cấu hình $\times$ 2 hệ phần cứng:

\begin{table}[h]\centering\small
\begin{tabular}{lccccc}\toprule
Ô & greedy & \texttt{exec3} & \texttt{llm3} & \texttt{exec3} phá & \texttt{llm3} phá\\\midrule
HE 1.5B (Kaggle) & 0,5375 & \textbf{0,6000} & 0,4812 & \textbf{0,0} & 2,8\\
HE 1.5B (5090)   & 0,5625 & \textbf{0,6438} & 0,4375 & \textbf{0,0} & 4,6\\
HE 7B (5090)     & 0,8000 & \textbf{0,8812} & 0,7812 & \textbf{0,0} & 2,6\\
HE 7B (Kaggle)   & 0,7938 & \textbf{0,9000} & 0,7438 & \textbf{0,0} & 3,2\\\bottomrule
\end{tabular}
\caption{Kiểm bằng chạy test (\texttt{exec3}) so với LLM tự kiểm (\texttt{llm3}).}
\end{table}

\texttt{exec3} phá 0 bài trong 20/20 fold và đạt đúng \texttt{oracle@4} ở hai ô có ghi
số; \texttt{llm3} phá ở cả 20 fold. Mốc trung thực là greedy --- bỏ phiếu có hại trên
code ($-0{,}113$/$-0{,}131$; lời giải dài hiếm trùng): $\texttt{exec3}-\text{greedy} =
+0{,}063 \to +0{,}106$. Giới hạn: khi model nền bão hoà (GSM8K 7B, solver 0,916), bộ kiểm thực
thi chính xác 0,837 vẫn cho hiệu ứng ròng âm (26 phá / 6 sửa).

\textbf{Tín hiệu học được} --- bộ phân loại huấn luyện trên lỗi tiêm (đổi một chữ số trong lời
giải đúng), cả hai ô hợp lệ:

\begin{table}[h]\centering\small
\begin{tabular}{lcccc}\toprule
Model & Lỗi tiêm & Lỗi thật & AUC & $\Delta$ độ chính xác khi dùng để chọn\\\midrule
1.5B & $-0{,}012$ & $+0{,}195$ & 0,528 & $-0{,}008$ (1/5 fold)\\
\textbf{7B} & $\mathbf{+0{,}573}$ & $\mathbf{+0{,}693}$ & \textbf{0,893} & $\mathbf{+0{,}024}$ (2/5 fold)\\\bottomrule
\end{tabular}
\caption{Bộ phân loại học được: AUC cao nhưng chuyển hoá thành thay đổi độ chính xác nhỏ.}
\end{table}

Ở 7B, AUC đạt 0,893 và hiệu năng trên lỗi thật cao hơn trên lỗi tiêm --- nhưng tín hiệu này
chỉ chuyển thành $+2{,}4$ điểm độ chính xác: tín hiệu đo được không tự động là tín hiệu dùng
được. Phát biểu hợp nhất, nối bốn thí nghiệm:

\begin{quote}
\textbf{Giá trị của một bộ kiểm bằng khoảng cách $\texttt{oracle@}k - \texttt{maj@}k$, không
bằng chất lượng bộ kiểm.} Code: khoảng cách $+21{,}3$ điểm --- bộ test lấy toàn bộ. Toán:
khoảng cách nhỏ --- AUC 0,893 lấy $+2{,}4$. Nút thắt nằm ở khâu sinh, không phải khâu
chọn (ở 50--58\% bài không đồng thuận, pool gần như không chứa ứng viên đúng).
\end{quote}

Hai trần định tuyến tính được trên dữ liệu \S\ref{sec:exposure}: cổng \textbf{theo artifact}
(chỉ cho xem khi artifact đúng --- oracle): 0,7160 so với $I$ 0,6980 = $\mathbf{+0{,}018}$; bộ
phân loại thật cần $\sim$89\% mới hoà vốn. Router \textbf{biết-khi-$I$-sai} (oracle mạnh hơn):
$+3{,}2$ điểm MATH --- nhưng đòi tín hiệu phát hiện \emph{lỗi của chính model mạnh}, đúng loại
tín hiệu mà thí nghiệm trên cho thấy đo được nhưng không chuyển thành độ chính xác. Hàm ý thực
tiễn: cải thiện bộ cổng có trần thấp ($+0{,}018$ với oracle); trong thiết lập này, cấu hình
mặc định không đưa artifact vào ngữ cảnh có kỳ vọng cao hơn.

% =====================================================================
\section{Tổng hợp: nghịch lý và lời giải}\label{sec:synthesis}

\textbf{Nghịch lý.} \S\ref{sec:roles}: chênh lệch năng lực lớn $\Rightarrow +14{,}0$ điểm.
\S\ref{sec:termA}: chênh lệch năng lực lớn $\Rightarrow \dceil$ âm. Cùng một biến, hai dấu.

\textbf{Lời giải} --- hai giao thức, một đẳng thức:
\begin{align*}
\text{Tuyển chọn (verifier, } 43{:}1\text{):} &\quad \text{value} \approx A \times \kappa - 0
  &&\to\ \text{chênh tăng} \Rightarrow \text{giá trị \textbf{tăng}}\\
\text{Sửa chữa (}V\text{ phải trả }B\text{):} &\quad \text{value} = A - B + C
  &&\to\ B \text{ tăng nhanh hơn } A \Rightarrow \text{giá trị \textbf{giảm}}
\end{align*}
Chênh lệch năng lực nuôi \textbf{cả $A$ lẫn $B$}; giao thức quyết định gặt bên nào. Năm mảnh
bằng chứng nội bộ hội tụ: (1) đại số --- chỉ-chọn đặt $B \approx 0$ theo cấu trúc; (2) độ lớn
--- trên MATH ở chênh lớn, $B$ gấp $\sim$11 lần $A$; (3) cơ chế --- $B$ là thiệt hại do tiếp
xúc nội dung sai, tiền đăng ký, hai miền; (4) trần --- cổng lọc hoàn hảo chỉ thu $+0{,}018$
trên nền $-0{,}124$; (5) liên miền --- quy luật MBPP dự báo được MATH ở 2/3 cặp.
Một quan sát ngoài --- tranh biện kém self-consistency ở model nhỏ trong số liệu công bố
\todoD{nguồn} --- nhất quán về hướng nhưng chờ xác nhận nguồn.

Bức tranh cuối là một \textbf{cây quyết định}:
\begin{itemize}\itemsep2pt
\item \textbf{Miền có bộ kiểm đúng đắn} (code chạy test): phối hợp đạt đúng trần oracle, 0 phá
      --- dùng, và dùng dạng chỉ-chọn.
\item \textbf{Miền không có:} bỏ phiếu cơ học lấy phần lớn giá trị của việc sinh nhiều lượt;
      thành phần LLM tổng hợp trung tính-đến-hại; tín hiệu học được đo được nhưng không dùng được.
\item \textbf{Trong các miền đã khảo sát:} cấu hình cho model mạnh xem bài của model yếu để
      sửa cho hiệu ứng âm hoặc không dương, rõ nhất khi chênh lệch vượt $\sim$0,09; gọi thẳng
      model mạnh là mốc mà không cấu hình sửa nào đã đo vượt được.
\end{itemize}

\begin{quote}\centering
Tóm lại: trong phạm vi khảo sát, cấu hình sinh độc lập rồi tuyển chọn\\
cho kết quả cao hơn các cấu hình cho model sửa bài của nhau.
\end{quote}

% =====================================================================
\section{Phương pháp luận khảo cứu}\label{sec:method}

Nhóm dùng hai chuẩn kiểm chứng bổ sung nhau: \textbf{thanh sai số theo fold} (khối vai trò/chi
phí) và \textbf{tiền đăng ký + điều kiện hợp lệ + niêm phong hash} (khối luồng thông tin: bảng
diễn giải khoá \emph{trước} khi chạy, có dòng bác bỏ giả thuyết; hash artifact ghi trước khi đọc
số; VOID thì không đọc). Các số liệu chính:
\begin{itemize}\itemsep2pt
\item \textbf{16/32 lần chạy có tệp kết quả mang trạng thái VOID (50\%)} --- hệ điều kiện hợp lệ
      đang làm đúng việc; 18 lần có niêm phong hash.
\item \textbf{Sổ dự đoán trước công khai: 21 đúng / 43} --- người đặt giả thuyết đoán sai một
      nửa; đó là lý do phải khoá diễn giải trước khi thấy số.
\item \textbf{Tính tất định:} hai tài khoản, hai ngày, cùng phần cứng $\Rightarrow$ giống nhau
      499/499 bài. Chạy lại cùng cấu hình không phải bằng chứng độc lập.
\item \textbf{Rà lại thống kê giữa kỳ:} ngưỡng ``5 điểm'' gốc suy cho phép đo đơn lẻ; với 5
      fold, chuẩn đúng là t-test (ngưỡng hiệu dụng $\sim$3,3 điểm) kèm KTC và xác suất một phía
      $P(\text{hiệu ứng} > 0)$. Rà 131 đại lượng: bốn kết quả chủ lực ($+14{,}0$, $+11{,}0$, $+5{,}6$,
      $+4{,}4$) giữ ý nghĩa ($t = 3{,}3$--$6{,}9$; $P(>0) = 98{,}5$--$99{,}9\%$); 6 kết quả được phục hồi \emph{sau
      khi đối chiếu cờ hợp lệ} (thống kê không cứu được lần chạy vô hiệu); 1 kết quả mất ý
      nghĩa ($k=3$ fold, $t_{\mathrm{crit}} = 4{,}303$); 2 ca treo được phán dứt khoát bằng dữ
      liệu fold còn trên Kaggle --- không cần chạy lại. 131 phép thử $\Rightarrow$ kỳ vọng
      $\sim$6,6 dương tính giả \cite{bh1995}; các ca $p \in [0{,}02;\,0{,}05]$ đọc với dè dặt
      tương ứng.
\end{itemize}
Phần lớn cải thiện ghi nhận ban đầu không giữ được ý nghĩa thống kê sau kiểm chứng; chúng
tôi báo cáo tỷ lệ này như một quan sát về quy trình đo.

% =====================================================================
\section{Hạn chế}\label{sec:limits}

\begin{enumerate}\itemsep2pt
\item \textbf{$\dhonest$ chưa có kết luận.} Đại lượng triển khai được thật đã qua năm lần chạy
      đều hỏng vì giới hạn 14,6 GB của GPU miễn phí; phần lớn kết luận định lượng dựa trên chặn
      trên $\dceil$ --- và mọi chặn trên của dòng sửa đều âm hoặc không dương có ý nghĩa.
\item \textbf{Mâu thuẫn chưa giải:} hai lần chạy 5 fold của \texttt{A\_gain} trên MATH 1.5B cho
      dấu ngược nhau ($-0{,}064$ và $+0{,}080$; khác cỡ fold 100/40, có thể khác xử lý
      \verb|\boxed|). Con số $-6{,}4$ chỉ trích kèm chú thích này cho tới khi giải xong.
\item Hai khối dùng hai chuẩn kiểm chứng, chưa kiểm chéo; đề xuất hợp nhất chuẩn chưa được nhóm
      phê duyệt.
\item Tính chuyển miền dựa trên 3 cặp, KTC rộng, 1/3 lệch hệ thống; vùng dương của quy luật
      chênh \textbf{không thể} xác lập trên MBPP (thiếu lực $\sim$8 lần).
\item $\kappa$ chưa giải được ở miền không có bộ kiểm đúng đắn: ba tín hiệu định tuyến đã thử
      đều thất bại hoặc trần quá thấp ($+0{,}018$ / $+3{,}2$ đều là oracle).
\item Khối luồng thông tin chỉ dùng greedy --- mỗi cấu hình một điểm; khối vai trò có fold bù lại.
\item Tập model bị giới hạn bởi VRAM tầng miễn phí; phạm vi khẳng định: \textbf{model
      nhỏ-đến-trung (1,5B--32B) trên bài toán suy luận}; chưa nói được về model rất lớn.
\item \S\ref{sec:router} và các phân rã MBPP ở \S\ref{sec:exposure} là mức B (đo một lần / hậu
      nghiệm); token 1.5B trong hiệu chỉnh FLOP là ước lượng từ ký tự (độ nhạy đã quét).
\end{enumerate}

% =====================================================================
\section{Kết luận}

``Đa tác tử có giúp không'' không có câu trả lời duy nhất; câu trả lời phụ thuộc ba thừa số
đo được. Dưới ba phép kiểm soát, giá trị của phối hợp phân rã thành ba thừa số đo được: dư địa
$H$ do chênh lệch năng lực quyết định; bộ chọn $\kappa$ do tính đúng đắn của tín hiệu quyết
định; thiệt hại $D$ do tiếp xúc với nội dung sai quyết định. Cùng một chênh lệch năng lực nuôi
cả cơ hội lẫn thiệt hại --- giao thức tuyển chọn gặt cơ hội, giao thức sửa chữa gánh thiệt hại;
và giao thức sửa, theo số liệu \S\ref{sec:exposure}, hoạt động như một ống dẫn đáp án hơn là
một bộ sửa lỗi. Nơi có bộ kiểm đúng đắn, phối hợp chạm trần oracle mà không phá bài nào; nơi
không có, cấu hình cho kết quả cao nhất trong khảo sát là sinh độc lập và bỏ phiếu, không đưa
bài làm của model này vào ngữ cảnh của model khác.

% =====================================================================
\begin{thebibliography}{99}\small
\bibitem{austin2021} J. Austin et al. \emph{Program Synthesis with Large Language Models.}
arXiv:2108.07732, 2021.
\bibitem{bh1995} Y. Benjamini, Y. Hochberg. \emph{Controlling the False Discovery Rate.}
JRSS--B 57(1), 1995.
\bibitem{chen2021} M. Chen et al. \emph{Evaluating Large Language Models Trained on Code.}
arXiv:2107.03374, 2021.
\bibitem{cobbe2021} K. Cobbe et al. \emph{Training Verifiers to Solve Math Word Problems.}
arXiv:2110.14168, 2021.
\bibitem{du2024} Y. Du et al. \emph{Improving Factuality and Reasoning in Language Models
through Multiagent Debate.} ICML, 2024.
\bibitem{gou2024} Z. Gou et al. \emph{CRITIC: Large Language Models Can Self-Correct with
Tool-Interactive Critiquing.} ICLR, 2024.
\bibitem{grattafiori2024} A. Grattafiori et al. \emph{The Llama 3 Herd of Models.}
arXiv:2407.21783, 2024.
\bibitem{guo2024} D. Guo et al. \emph{DeepSeek-Coder: When the Large Language Model Meets
Programming.} arXiv:2401.14196, 2024.
\bibitem{hendrycks2021} D. Hendrycks et al. \emph{Measuring Mathematical Problem Solving with
the MATH Dataset.} NeurIPS, 2021.
\bibitem{huang2024} J. Huang et al. \emph{Large Language Models Cannot Self-Correct Reasoning
Yet.} ICLR, 2024.
\bibitem{kaplan2020} J. Kaplan et al. \emph{Scaling Laws for Neural Language Models.}
arXiv:2001.08361, 2020.
\bibitem{lightman2024} H. Lightman et al. \emph{Let's Verify Step by Step.} ICLR, 2024.
\bibitem{madaan2023} A. Madaan et al. \emph{Self-Refine: Iterative Refinement with
Self-Feedback.} NeurIPS, 2023.
\bibitem{shapley1953} L. S. Shapley. \emph{A Value for $n$-Person Games.} Contributions to the
Theory of Games II, 1953.
\bibitem{shinn2023} N. Shinn et al. \emph{Reflexion: Language Agents with Verbal Reinforcement
Learning.} NeurIPS, 2023.
\bibitem{snell2024} C. Snell et al. \emph{Scaling LLM Test-Time Compute Optimally Can Be More
Effective than Scaling Model Parameters.} arXiv:2408.03314, 2024.
\bibitem{wang2023} X. Wang et al. \emph{Self-Consistency Improves Chain of Thought Reasoning in
Language Models.} ICLR, 2023.
\bibitem{wei2022} J. Wei et al. \emph{Chain-of-Thought Prompting Elicits Reasoning in Large
Language Models.} NeurIPS, 2022.
\bibitem{yang2024} A. Yang et al. \emph{Qwen2.5 Technical Report.} arXiv:2412.15115, 2024.
\bibitem{zheng2023} L. Zheng et al. \emph{Judging LLM-as-a-Judge with MT-Bench and Chatbot
Arena.} NeurIPS, 2023.
\end{thebibliography}

% =====================================================================
\appendix
\section*{Phụ lục (con trỏ nguồn)}
\begin{itemize}\itemsep1pt
\item[\textbf{A.}] Thuật ngữ và ký hiệu --- \texttt{report/THUAT\_NGU.md} (chép mục 1--4 khi dựng bản in).
\item[\textbf{B.}] Trích các bản tiền đăng ký (\#104, \#106, \#107, \#112) --- \texttt{docs/PREREGISTRATION.md}.
\item[\textbf{C.}] Bảng niêm phong hash --- \texttt{docs/RESULT\_SEALS.md}.
\item[\textbf{D.}] 16 lần chạy VOID (trên 32) và lý do.
\item[\textbf{E.}] Sổ dự đoán trước (21/43).
\item[\textbf{F.}] 37 quy tắc quy trình --- \texttt{docs/QUY\_TRINH\_VONG\_LAP.md}.
\item[\textbf{G.}] Bảng kết quả đầy đủ khối nhóm --- \texttt{docs/RESULTS.md}.
\item[\textbf{H.}] Phân tích Shapley theo vai --- 29 tài liệu trong \texttt{docs/}.
\item[\textbf{I.}] Rà lại thống kê giữa kỳ và các phán quyết --- \texttt{CAU\_HOI\_THAO\_LUAN.md} nhóm E, mục A6--A7.
\end{itemize}

\paragraph{Hình dự kiến.} (1) Sơ đồ khung $H \times \kappa - D$ + ánh xạ ba nhánh ---
\emph{Nguyên}; (2) \textbf{nghịch lý hai đường ngược chiều} (quan trọng nhất) --- \emph{Nguyên};
(3) $\dceil$ theo chênh lệch, 15 điểm + $g^\ast$ --- \emph{Tùng Dương}; (4) phân tầng tiếp xúc
$2\times2$ --- \emph{Tùng Dương}; (5) accuracy--chi phí (trọng số FLOP), có điểm router ---
\emph{Tùng Dương}; (6) bất đối xứng verifier, có KTC --- \emph{Đức}; (7) trần định tuyến theo độ
chính xác bộ phân loại --- \emph{Tùng Dương}.

\end{document}
